\section{Resultants and divisor multiplicities}

Because $m$ is monic, taking the norm from $\Q(\theta)$ is equivalent to a
resultant in $\theta$.  Exact elimination gives
\[
 \Res_{\theta}(m,q_0)=P_{18}/49,
 \qquad
 \Res_{\theta}(m,q_1)=P_{12}/7.
\]
Consequently
\begin{equation}\label{eq:norm}
 N_7(z)=\frac{P_{18}(z)^2}{49P_{12}(z)^2},
 \qquad N_7(0)=1.
\end{equation}

The numerator and denominator of \eqref{eq:norm} have degrees $36$ and $24$.
Their difference $12$ agrees with the general ordinary-norm formula
$2(p-1)$.

\begin{proposition}[Exact gcd certificate]\label{prop:gcd}
The polynomials $P_{18}$ and $P_{12}$ are squarefree and mutually coprime over
$\Q$.
\end{proposition}

\begin{proof}
Reduction modulo five preserves both degrees.  Direct Euclidean algorithms in
$\F_5[z]$ give
\[
 \gcd(P_{18},P_{12})=gcd(P_{18},P_{18}')
 =\gcd(P_{12},P_{12}')=1.
\]
A common factor or repeated factor over $\Q$ would give one after reduction
at a degree-preserving good prime.  Thus the three rational gcds are one.
\end{proof}

Every zero of $N_7$ now has valuation two and every pole valuation minus two.
The points lie on $|z|=1$: before Galois descent the factors are characteristic
polynomials of unitary sector blocks, and \cref{prop:gcd} excludes cross-sector
cancellation.

